\chapter{30}


\emph{«Tagwol»}\footnote{Gängiges Grusswort der einheimischen Walliser.}.

Der Mann begrüsste ihn zuerst.

\emph{«Tagwol»}, rief ihm nun auch die Frau zu.

Sie lösten ihre Umarmung und musterten ihn.

Louis blieb einige Schritte vor ihnen stehen.

Er sprach mit fester Stimme.

«Guten Tag, bin ich hier richtig, \emph{Sunnewaid 2}?»

«Ja, willkommen. Was führt dich zu uns?»

Der Mann strahlte über das ganze Gesicht. Auch die Frau schien sich über
Louis’ Anblick zu freuen. Sie griff sich einen Apfel aus einer Schale
auf dem Tisch, ging ein paar Schritte auf Louis zu und lächelte
freundlich.

«Ich bin Ciro und habe euer Inserat gelesen. Ihr sucht einen
\emph{Zusenn?»}

«Ich fass es nicht, Joh,» jubelte die Frau mit einem Seitenblick zum
Mann, «wir wurden doch noch erhört.»

Sie stellte sich vor Louis und reichte ihm die Hand.

«Manuela, hi, wie, wie heisst du nochmal?»

«Ciro, Ciro Lorente.»

Louis erwiderte ihren Händedruck.

«Hallo, Manuela.»

Der Mann klopfte Louis auf die Schulter und lachte nun herzlich.

«Dich schickt uns der Himmel, Ciro. Bin Johannes, kurz Joh. Willkommen
auf der \emph{Sunnewaid}. Stell doch dein Gepäck da an die Mauer und setz dich. Du siehst durstig aus.»

Er führte Louis zur Bank vor dem schweren Holztisch neben dem
Hauseingang und rückte ihm einen Stuhl zurecht.
Manuela goss ihm Apfelsaft in ein Glas und Louis fühlte sich willkommen wie ein verlorenes Kind, das nach
erfolgloser Suche seiner Eltern, eines Tages unerwartet zurückgekehrt
war. Während er trank, schauten sie ihn
ununterbrochen an. Direkt und ungemein wach. Louis hielt sich an seinem
Glas fest. Joh lehnte sich zurück. Manuela atmete ein, während sie ihr
langes Haar zurückwarf.

«Erzähl, Ciro, wo kommst du her? Wie hast du von uns erfahren?»

Sie hatte eine wohlklingende Stimme. In ihrer Präsenz strahlte sie eine
Souveränität aus, die Louis gefiel und ihn gleichzeitig einschüchterte.
Nach wie vor hielt sie Augenkontakt. Auch Joh, der sein Kinn etwas
vorstreckte, fixierte Louis. Permanent. Seine dunklen Augen wirkten
durch die Brillengläser vergrössert. Bergbraungebrannt und mit seinen
markanten Gesichtszügen erinnerte er Louis entfernt an Al Pacino ohne
Bart.

Louis erzählte mit wenigen Worten von einer angeblichen Auszeit. Er finanziere seine Reise vorwiegend mit Strassenmusik und es sei Zeit, etwas dazuzuverdienen.

« - und dann hab ich einen Typen in Zürich kennengelernt, der mir
erzählte, er sei auf dem Weg auf eine Alp, den Namen habe ich vergessen - 
er brauche ein paar Wochen Ruhe und werde sich auf einem Hof von
Bekannten nützlich machen. Ich wurde hellhörig. Ich sagte ihm, das
klinge gut, das würde mich auch interessieren. Und dann kam das eine zum
andern und er zeigte mir einige Links zu Bergbauern, die für den Sommer und den
kommenden Herbst nach Verstärkung, nach Saisonarbeitern, suchten, nach Schafhirten auch. So
bin ich auf euren Aufruf der Suche nach einem \emph{Zusenn} gestossen.
Und jetzt bin ich hier.»

Die beiden hatten ihre konzentrierten Blicke keinen Moment von ihm
abgewandt. Irritiert wirkten sie nicht. Stumm hatten sie Louis zugehört,
zwischendurch anteilnehmend genickt und schliesslich erfreut ausgesehen.
Sie schienen weitere Erklärungen zu erwarten.

«Ich, ich weiss nicht, was ihr genau sucht, was es zu tun gibt bei euch.
Und ehrlich gesagt, bin ich vollkommen unwissend, aber bereit, was Neues
zu lernen.»

Es folgte ein Moment des Wartens. Louis wurde nervös. Es war
ungewohnt für ihn, nicht unterbrochen zu werden. Erst als
Joh mit funkelnden Augen einatmete, fühlte er sich erlöst. Der Bann
schien gebrochen.

«Klingt gut, Ciro,» begann er und rückte seine Brille zurecht, die ihm
unterdessen der Nase entlang nach unten gerutscht war, «und wie gesagt,
du könntest unsere Rettung sein. Wir brauchen dringend einen zweiten
Hirten auf den Alp-Höhen oben, um den Aletschgletscher, wir haben dort
unsere Schafe. Und Mitte September nach der \emph{Schafscheid} auf der
Bettmeralp haben wir hier unten Hochbetrieb beim Schafe scheren. Du
siehst stark aus, widerstandsfähig, der Rest ist lernbar. Willkommen
schon mal bei uns.»

Damit stand Joh zackig auf. Er ging um den Tisch herum auf Louis zu und
drückte ihm kräftig die Hand.

«Und mutig bist du offensichtlich auch, einfach hier rauf zu kommen ohne
vorher anzurufen. Das gefällt mir.»

Er klopfte Louis auf die Schulter und entschuldigte sich, heute sei
Manuelas Monats-Tag. Er habe zu tun.

«Schauen wir bei Apéro und Nachtessen, ob wir uns finden, ok?»

Louis nickte. Die Aussicht auf Essen klang vielversprechend.

«Komm erstmal an,» meinte Joh weiter und nach einem klärenden Blick zu
Manuela, «fürs Erste kannst du heute gerne mal bei uns übernachten.
Unser Hof ist unter anderem eine Herberge für gelegentliche Gäste.»

Wieder nickte Louis. Diesmal ungläubig. Er griff nach dem Glas und trank
einen Schluck.

«Genau,» Manuela wies Louis mit einer Handbewegung den Weg, «wenn du
möchtest, zeige ich dir das Gastzimmer, ok?»

«Sehr gerne, danke, das ist ein tolles Angebot.»

Joh verschwand beschwingt die wenigen Treppenstufen hinauf ins Haus.
Louis blickte ihm nach. Die waren schnell unterwegs hier oben. Er staunte.
Zusammen mit Manuela betrat Louis das Haus über den Kücheneingang. Sie
führte ihn in einen anschliessenden Raum, von dem eine enge Treppe steil
hinauf in einen schmalen Flur führte. Vorsichtig lenkte Louis Rucksack
und Gitarrenkoffer an den Wänden vorbei. Manuela wies mit der Hand dem
Flur entlang nach rechts.

«Schau, du kannst im \emph{Schäferloft} übernachten.»

Sie machte den Eindruck, als freue sie sich beim Aussprechen des
Zimmernamens wie ein Kind.
Die Ausstattung des Raumes überraschte Louis.

«Wow, was für ein schönes \emph{Loft}.»

Er lachte.
Die Zimmertür war mit schwungvollen, dunkelblauen Buchstaben
beschriftet. Wie alle Zimmer hier oben, wie er später bemerkte.

«Komm einfach runter, wenn du soweit bist. Nimm dir Zeit.»

Manuela lächelte Louis zu und liess ihn alleine. Er hörte sie die
knatternden Holzstufen nach unten steigen und schloss die Tür hinter
sich.
Mit Rucksack und Gitarrenkoffer in den Händen stand Louis inmitten des
Zimmers. Es zeigte sich holzgetäfert, hell gestrichen mit zwei
Einzelbetten. Eines an der Innenwand, das andere unter dem Fenster.
Weisse Vorhänge vor der aussergewöhnlich breiten Fensterfront fielen
luftig zu Boden. Der weite Stoff liess das Gemach in artfremdem Luxus
erscheinen. Der Raum duftete, als wäre er gerade eben einer Waschmaschine 
entstiegen. Louis schmunzelte, atmete tief ein, stellte Rucksack und Gitarrenkoffer ab 
und liess sich auf das Bett unter dem Fenster fallen. Er seufzte. Anspannung und
Unsicherheit schienen auf der Stelle von ihm abzufallen. Sein
anfängliches Staunen blieb. Der Ort war einzigartig, anders als alles, was
er kannte. Louis sah an die Decke und suchte nach einem passenden
Begriff, um die Stimmung festzumachen. Gemütlich, dachte er, war nicht
treffend genug. 

Plötzlich riss ihn ein lauter Knall aus dem Untergeschoss aus seinen
Gedanken. Louis setzte sich schlagartig auf. Es hatte sich angehört, als
wäre eine Tür zugeschlagen und Glas zerschmettert worden. Darauf hörte
er Joh laut und schnell sprechen. Er klang aufgebracht schnittig.

« - niemals mehr. Du kennst die Bedingungen, du hast zugestimmt und
hältst dich daran. Das gilt für beide Seiten.»

So viel zu \emph{gemütlich,} schoss Louis durch den Kopf. Was mochte da
unten los sein?
Eine ihm unbekannte Männerstimme zwängte sich zwischen Johs Worte. Nun
sprachen sie leiser. Und Louis verstand nur mehr einzelne Wortfetzen. Er
stand leise auf und schlich sich unter gelegentlichem Bodenknarzen an
die Tür. Er öffnete sie einen Spalt weit, hielt den Atem an und
lauschte.

«Das geht dich nichts an. Kümmere dich um dein eigenes Zeug,» die fremde
Stimme explodierte, «ich kann damit umgehen. So viel zu Verantwortung,
verdammt.»

Joh war nun ruhiger und leiser zu hören.

«Null Toleranz ist null Toleranz. Ich weiss nicht, was du nicht
verstehst. Es gibt nichts zu verhandeln, David. Du kennst die Regeln.
Pack deine Sachen. Um achtzehn dreissig fährt die nächste Bahn runter.
Nächsten Freitag bist du wieder willkommen. Ich sorge dafür, dass du
unten abgeholt wirst.»

Wieder knallte es. Diesmal gedämpfter. Womöglich war ein Stuhl
umgekippt. Schritte krachten schwer über den Küchenboden und drangen nun
die Treppe hoch. Louis schloss erschrocken die Zimmertür. Er schwitzte.
Noch war es möglich, den Ort wieder zu verlassen. Nachdenklich zog er sich die Schuhe aus.
Bevor er nach unten ging, wollte er sich wenigstens Gesicht und Hände waschen. Eine Katzenwäsche musste reichen. 
Was würde ihn wohl unten erwarten?

Seine Hygieneartikel in den Händen stellte er sich vor den Spültisch.
Dabei schaute er direkt in einen länglichen, abgerundeten Spiegel.
Er zuckte zusammen. Ein Fremder starrte ihn an. Ein Wilder mit einem
abgekämpft müden Blick, einem jung spriessenden Bart und Haaren, die in
alle Richtungen standen. Doch je länger er sein Gesicht anstaunte, desto
zufriedener stimmte ihn das Bild. Ciro, das war Ciro. Was konnte besser
sein, als dass er sich selbst kaum mehr erkannte.
Nachdem er sich zurecht gemacht und sein neues T-Shirt übergezogen
hatte, richtete er sich auf. Der Auftritt konnte kommen. Bereits im Flur
stieg Louis ein Duftgemisch aus Kräutern, gebratenen Kartoffeln und
Fleisch in die Nase. Sein Magen begann verrückt zu spielen. Er verspürte
einen masslosen Hunger. Aus dem letzten Zimmer am Ende des Flurs vernahm
er Geräusche und ein Fluchen. Die Tür war oben mit \emph{Boxenstopp}
angeschrieben.
Louis stieg die Treppe hinunter in die grosse Küche.
Joh schenkte gerade in drei hohe Gläser eine rote Flüssigkeit ein. Die
Campari-Flasche stellte er zurück auf ein gläsernes, langes Regal neben
der Abzugshaube. Louis liess seinen Blick interessiert an den dicht
nebeneinander aufgereihten Spirituosen entlang gleiten. Manuela sah er
durch das geöffnete Fenster draussen auf der Bank sitzen. An die
vorangegangene Szene erinnerte nur eine fehlende Scheibe im oberen Teil
der Eingangstür zur Küche. Joh gab sich ebenso entspannt wie bei Louis’
Ankunft.

«Ciro. Schön, bist du bereit,» er klang fröhlich, «magst du einen
Aperitif? Ich habe dazu Bruschetta vorbereitet.»

Er zeigte beschwingt auf eine Glasplatte mit einladend arrangierten
Tomatenbrötchen.

«Ja, gerne.»

Louis schaute Joh verwundert nach, wie der nun auf eine Musikanlage
zusteuerte und die Fernbedienung holte. Dann wies er Louis mit einer Handbewegung zwinkernd nach draussen, drückte auf Play, schnappte sich die Platte und das Tablett mit den Gläsern und stieg die wenigen Treppenstufen hinter Louis auf den Vorplatz hinunter.

Manuela sah anders aus als zu Louis’ Ankunft. Das eng anliegende, 
chice Sommerkleid betonte ihre Figur. Ihr langes Haar hatte sie
hochgesteckt. Sie hatte sich offensichtlich feingemacht. Sie stand auf
und schien geheimnisvoll zu leuchten. Und schon begannen sich gut hörbar erste Klänge eines Musikstücks aus
der Küche hinaus auf den Platz vor dem Haus zu legen. Nach dem Geräusch
eines Flugzeugmotors und den Worten einer weiblichen Stimme setzte ein
Saxophon ein. Männlicher Gesang folgte.

\qrblock{Taxi \emph{«Campari Soda»}}{media/QR-Codes/027_Campari_Soda_Taxi_Spotify.png}

Die Stimme hörte sich abgeklärt an. Louis kannte den Song nicht. Das
Lied klang leicht, einfach und zum Abheben schön. Im wahrsten Sinne. Wie
zuvor, als er die beiden aus der Ferne beobachtet hatte, lag Magie in 
der Luft. 
Joh reichte Louis ein Glas. Dann wandte er sich an Manuela.
Er zauberte sich ein verführerisches Lächeln ins Gesicht, legte den Kopf schräg und blickte sie charmant an.

«Campari Soda für meine Lady und dazu wird natürlich der Blick aufs Wolkenmeer serviert.»

Mit der Galanterie eines Kellners erster Klasse verneigte er sich. Und
jetzt fiel Louis auf, dass auch Joh sich umgezogen hatte. Er trug ein
weisses Hemd, exakt gebügelt, das lose über seine Jeans hing. Manuela
nickte erwartungsvoll, während sie den Cocktail entgegen nahm. Joh legte
seine freie Hand über ihre Wange und küsste sie auf die Stirn.
Louis beobachtete die beiden amüsiert. Filmszene, 2. Akt.
Während sie sich setzten, fiel Louis ein viertes Gedeck auf, das Joh mit
wenigen geschickten Handbewegungen auf einen Stuhl am Tischende verschwinden
liess.
